\section{Introduction}
\label{introduction}

Recent studies have increasingly emphasized the central role of requirements in AI-driven software engineering.
Robinson et al.~\cite{robinsonRequirementsAreAll2025} envision a future paradigm in which end-users can directly drive the entire software development lifecycle through multimodal requirements expressed in natural language, images, audio, and video, highlighting requirements as the core interface between humans and generative software systems.
Along this direction, Han et al.~\cite{hanArchCodeIncorporatingSoftware2024} proposed ARCHCODE, a requirement-aware code generation framework that explicitly extracts and extrapolates both functional and non-functional requirements from textual descriptions to guide code and test generation, demonstrating that requirement-centric reasoning can substantially improve code correctness and requirement satisfaction.
Meanwhile, Darif et al.~\cite{darifControlledNaturalLanguage2026} systematically reviewed Controlled Natural Language (CNL) approaches for requirements specification and showed that reducing ambiguity while preserving stakeholder readability remains a longstanding challenge in requirements engineering.
Their findings further reveal limitations in automated tooling, domain vocabulary integration, and practical adoption, suggesting that intelligent and cognitively guided requirement understanding remains an open research problem.
Together, these studies indicate a growing shift from syntax-oriented software generation toward requirement-centric and cognition-aware software engineering paradigms.